\section{Expresión matemática de RoPE}

\subsection{Propiedades básicas de la matriz de rotación}

Se aplica la matriz de rotación $R_m$ al Query de la posición $m$, y $R_n$ al Key de la posición $n$:

$$
\text{score}(m, n) = (R_m q)^T (R_n k) = q^T R_m^T R_n k = q^T R_{n-m} k
$$

Propiedad clave: $R_m^T R_n = R_{n-m}$, la transpuesta de la matriz de rotación es igual a la rotación inversa, y la combinación de dos rotaciones solo depende de la diferencia de posiciones $n - m$.

\begin{warningbox}{La relatividad se deduce, no se construye}
Solo necesitamos calcular de antemano la matriz de rotación para cada posición absoluta; la codificación de posición relativa es un resultado natural de la deducción matemática y no requiere construirse de forma explícita.
\end{warningbox}

\subsection{Parámetro de frecuencia $\theta$}

RoPE aplica rotaciones de distinta frecuencia sobre los diferentes subespacios 2D (subspace) del vector de embedding:

$$
\theta_i = \text{base}^{-2i/d}
$$

\begin{itemize}
  \item $\text{base}$: generalmente se toma 10,000 (RoPE estándar); Qwen3 toma \textbf{1,000,000} (para lograr una ventana de contexto mayor)
  \item $i$: índice del subespacio, de 0 a $d/2 - 1$
  \item $d$: head dimension
\end{itemize}

$\theta_i$ decae de forma \textbf{exponencial} conforme aumenta $i$; al tomar $\log$, decae linealmente.

\subsection{Resumen de la sección}

RoPE logra la codificación de posición relativa mediante matrices de rotación de posiciones absolutas; $\theta_i$ controla la frecuencia de rotación de cada subespacio. Cuanto mayor sea base, más lenta será la rotación de los subespacios de baja frecuencia, lo cual favorece el modelado de largo alcance.
